% ===========================================================================
%  Annexe — Formulaire
%  Conçu pour être détaché et gardé sous les yeux toute l'année.
% ===========================================================================
\chapter{Formulaire}

Tout le programme de l'année tient dans les pages qui suivent. Elles sont
faites pour être détachées, pliées en deux et gardées dans le cahier — ou
recopiées à la main, ce qui vaut encore mieux.

% ---------------------------------------------------------------------------
\section{Second degré}

Pour $P(x) = ax^{2}+bx+c$ avec $a \neq 0$, et $\Delta = b^{2}-4ac$ :

\begin{center}
\renewcommand{\arraystretch}{1.4}
\begin{tabular}{|L{2.3cm}|L{4.9cm}|L{4.6cm}|}
\hline
\textbf{Si} & \textbf{Racines} & \textbf{Factorisation}\\
\hline
$\Delta>0$ & $x_{1,2} = \dfrac{-b \mp \sqrt\Delta}{2a}$ &
$a\left(x-x_{1}\right)\left(x-x_{2}\right)$\\
$\Delta = 0$ & $x_{0} = -\dfrac{b}{2a}$ (double) &
$a\left(x-x_{0}\right)^{2}$\\
$\Delta<0$ & aucune racine réelle & pas de factorisation dans $\R$\\
\hline
\end{tabular}
\end{center}

\[
  S = x_{1}+x_{2} = -\frac ba, \qquad P = x_{1}x_{2} = \frac ca .
\]

\textbf{Signe :} le trinôme est du signe de $a$ partout, sauf entre les
racines lorsqu'elles existent.

\textbf{Bicarrée :} $ax^{4}+bx^{2}+c = 0$ se résout en posant $X = x^{2}$,
avec $X \geq 0$.

% ---------------------------------------------------------------------------
\section{Systèmes linéaires}

\[
  \begin{cases} ax+by = c\\ a'x+b'y = c'\end{cases}
  \qquad
  D = \begin{vmatrix} a & b\\ a' & b'\end{vmatrix} = ab'-a'b .
\]
Si $D \neq 0$ : $x = \dfrac{D_{x}}{D}$ et $y = \dfrac{D_{y}}{D}$, où
$D_{x}$ et $D_{y}$ s'obtiennent en remplaçant la colonne correspondante par
celle des seconds membres. Si $D = 0$ : aucune solution, ou une infinité.

\textbf{Ordre 3, règle de Sarrus} — on recopie les deux premières colonnes à
droite, on ajoute les trois produits descendants et l'on retranche les trois
montants :
\[
  D = ab'c''+bc'a''+ca'b''-ca''b'-ab''c'-bc''a' .
\]

% ---------------------------------------------------------------------------
\section{Limites}

\[
  \limpinf \frac1x = 0, \quad
  \lim_{x\to0^{+}}\frac1x = +\infty, \quad
  \lim_{x\to0^{-}}\frac1x = -\infty, \quad
  \limpinf x^{n} = +\infty .
\]

\textbf{Formes indéterminées :} $\infty-\infty$, $0\times\infty$,
$\dfrac{\infty}{\infty}$, $\dfrac00$.

\textbf{À l'infini :} un polynôme se comporte comme son terme de plus haut
degré ; un quotient de polynômes, comme le quotient de leurs termes de plus
haut degré.

\textbf{En $\frac00$ :} factoriser par $(x-a)$ au numérateur et au
dénominateur, puis simplifier.

\begin{center}
\renewcommand{\arraystretch}{1.4}
\begin{tabular}{|L{5.4cm}|L{6.4cm}|}
\hline
\textbf{Limite} & \textbf{Interprétation graphique}\\
\hline
$\displaystyle\lim_{x\to a} f(x) = \pm\infty$ & asymptote verticale
$x = a$\\
$\limpinf f(x) = \ell$ & asymptote horizontale $y = \ell$\\
$\limpinf\left[f(x)-(ax+b)\right] = 0$ & asymptote oblique $y = ax+b$\\
$\limpinf f(x) = \pm\infty$ et $\limpinf\frac{f(x)}{x} = \pm\infty$ &
branche parabolique de direction $(Oy)$\\
\hline
\end{tabular}
\end{center}

\textbf{Position :} le signe de $f(x)-(ax+b)$ dit si la courbe est au-dessus
(positif) ou en dessous (négatif) de la droite.

% ---------------------------------------------------------------------------
\section{Dérivation}

\begin{center}
\renewcommand{\arraystretch}{1.5}
\begin{tabular}{|L{3.1cm}|L{3.1cm}||L{2.6cm}|L{2.6cm}|}
\hline
\textbf{$f$} & \textbf{$f'$} & \textbf{$f$} & \textbf{$f'$}\\
\hline
$k$ & $0$ & $u+v$ & $u'+v'$\\
$x^{n}$ & $n\,x^{n-1}$ & $ku$ & $k\,u'$\\
$\dfrac1x$ & $-\dfrac{1}{x^{2}}$ & $uv$ & $u'v+uv'$\\
$\sqrt x$ & $\dfrac{1}{2\sqrt x}$ & $\dfrac1v$ &
$-\dfrac{v'}{v^{2}}$\\
$\ln x$ & $\dfrac1x$ & $\dfrac uv$ & $\dfrac{u'v-uv'}{v^{2}}$\\
$e^{x}$ & $e^{x}$ & $u^{n}$ & $n\,u'\,u^{n-1}$\\
$\ln u$ & $\dfrac{u'}{u}$ & $e^{u}$ & $u'\,e^{u}$\\
\hline
\end{tabular}
\end{center}

\textbf{Tangente} au point d'abscisse $a$ :
$y = f'(a)\,(x-a)+f(a)$.

\textbf{Variations :} $f'>0$ sur $I$ $\Rightarrow$ $f$ croissante sur $I$ ;
$f'<0$ $\Rightarrow$ décroissante. Extremum là où $f'$ s'annule
\emph{en changeant de signe}.

% ---------------------------------------------------------------------------
\section{Logarithme népérien}

Défini sur $\intervalleo{0}{+\infty}$, strictement croissant,
$\ln 1 = 0$, $\ln e = 1$.

\[
  \ln(ab) = \ln a+\ln b, \qquad \ln\frac ab = \ln a-\ln b, \qquad
  \ln\left(a^{n}\right) = n\ln a, \qquad \ln\sqrt a = \frac12\ln a .
\]
\[
  \lim_{x\to0^{+}}\ln x = -\infty, \qquad \limpinf \ln x = +\infty, \qquad
  \limpinf \frac{\ln x}{x} = 0, \qquad \lim_{x\to0^{+}} x\ln x = 0 .
\]

\textbf{Équations :} $\ln u = a \daccord u = e^{a}$ ;
$\ln u = \ln v \daccord u = v$ ;
$a\ln^{2}x+b\ln x+c = 0$ se résout en posant $X = \ln x$.

\textbf{Toujours :} écrire les conditions d'existence avant de calculer, et
rejeter à la fin les solutions qui ne les vérifient pas.

% ---------------------------------------------------------------------------
\section{Exponentielle}

Définie sur $\R$, strictement croissante, strictement positive,
$e^{0} = 1$.

\[
  \ln\left(e^{x}\right) = x, \qquad e^{\ln y} = y \;(y>0),
\]
\[
  e^{a+b} = e^{a}e^{b}, \qquad e^{-a} = \frac{1}{e^{a}}, \qquad
  e^{a-b} = \frac{e^{a}}{e^{b}}, \qquad \left(e^{a}\right)^{n} = e^{na} .
\]
\[
  \limminf e^{x} = 0, \qquad \limpinf e^{x} = +\infty, \qquad
  \limpinf \frac{e^{x}}{x} = +\infty, \qquad \limminf x\,e^{x} = 0 .
\]

\textbf{Équations :} $e^{u} = a \daccord u = \ln a$ (et aucune solution si
$a \leq 0$) ; $e^{u} = e^{v} \daccord u = v$ ;
$a\,e^{2x}+b\,e^{x}+c = 0$ se résout en posant $X = e^{x}>0$.

% ---------------------------------------------------------------------------
\section{Primitives et intégrale}

$F$ est une primitive de $f$ sur $I$ lorsque $F' = f$ sur $I$. Deux
primitives diffèrent d'une constante.

\begin{center}
\renewcommand{\arraystretch}{1.5}
\begin{tabular}{|L{3.2cm}|L{3.2cm}||L{2.6cm}|L{2.6cm}|}
\hline
\textbf{$f$} & \textbf{une primitive} & \textbf{$f$} &
\textbf{une primitive}\\
\hline
$k$ & $kx$ & $\dfrac{1}{x^{2}}$ & $-\dfrac1x$\\
$x^{n}$ & $\dfrac{x^{n+1}}{n+1}$ & $\dfrac1x$ & $\ln x$ (sur $\Rps$)\\
$e^{x}$ & $e^{x}$ & $e^{ax}$ & $\dfrac1a e^{ax}$\\
$\dfrac{u'}{u}$ & $\ln u$ (où $u>0$) & $\ln x$ & $x\ln x-x$\\
\hline
\end{tabular}
\end{center}

\[
  \int_{a}^{b} f(x)\dx = \left[F(x)\right]_{a}^{b} = F(b)-F(a) .
\]

\textbf{Aire.} Si $f \geq 0$ sur $\intervalle{a}{b}$, l'aire du domaine
limité par la courbe, l'axe des abscisses et les droites $x = a$ et $x = b$
vaut $\int_{a}^{b} f(x)\dx$ unités d'aire ; si $f \leq 0$, c'est l'opposé.

\textbf{Conversion.} Unité graphique $u$ centimètres
$\Rightarrow$ $1$ unité d'aire $= u^{2}$ cm².

% ---------------------------------------------------------------------------
\section{Dénombrement et probabilité}

\begin{center}
\renewcommand{\arraystretch}{1.45}
\begin{tabular}{|L{4.9cm}|L{2.9cm}|L{3.4cm}|}
\hline
\textbf{L'énoncé dit} & \textbf{L'ordre} & \textbf{On compte avec}\\
\hline
successivement, avec remise & compte & $n^{p}$\\
successivement, sans remise & compte &
$\arrang{p}{n} = \dfrac{n!}{(n-p)!}$\\
simultanément, d'un seul coup & ne compte pas &
$\combi{p}{n} = \dfrac{n!}{p!(n-p)!}$\\
\hline
\end{tabular}
\end{center}

\[
  \combi{0}{n} = \combi{n}{n} = 1, \qquad \combi{1}{n} = n, \qquad
  \combi{p}{n} = \combi{n-p}{n}, \qquad 0! = 1 .
\]

\textbf{Équiprobabilité :}
$\proba{A} = \dfrac{\text{nombre de cas favorables}}
{\text{nombre de cas possibles}}$.

\[
  0 \leq \proba{A} \leq 1, \qquad \proba{\evtc{A}} = 1-\proba{A},
  \qquad \proba{A \cup B} = \proba{A}+\proba{B}-\proba{A \cap B} .
\]

\textbf{« Au moins un »} a pour contraire \textbf{« aucun »} : c'est
presque toujours le chemin le plus court.

\textbf{Hors programme en série L :} probabilités conditionnelles,
indépendance, variables aléatoires.

% ---------------------------------------------------------------------------
\section{Statistique à deux variables}

\textbf{Point moyen :} $\pmoyen\left(\moy x\,;\,\moy y\right)$ où $\moy x$
et $\moy y$ sont les moyennes des $x_{i}$ et des $y_{i}$.

\textbf{Méthode de Mayer :} on partage la série en deux sous-groupes de même
effectif, on calcule leurs points moyens $G_{1}$ et $G_{2}$, et la droite
d'ajustement est $\droiteMayer$. Son équation $y = ax+b$ s'obtient en
résolvant
\[
  \begin{cases} x_{G_{1}}\,a+b = y_{G_{1}}\\
                x_{G_{2}}\,a+b = y_{G_{2}} \end{cases}
\]
\textbf{Vérification gratuite :} $\pmoyen$ appartient à cette droite.

\textbf{Taux d'évolution :} $t = \dfrac{V_{1}-V_{0}}{V_{0}}$ ; coefficient
multiplicateur $1+t$. Deux évolutions successives se composent en
\emph{multipliant} les coefficients, jamais en additionnant les taux.
